% !TeX root = ../nc-del-pezzo.tex
\section{Poisson模型への代数的比較と相対幾何}
\label{sec:poisson-shadow-comparison}

定理\ref{thm:zero-neighborhood-finite-diagonal}の全可換零切断を含む近傍を用いる。
本節では，対称化した中心方程式に最低次数の交換子を組み合わせ，
可換模型上のPoisson構造を構成する。
その射を零切断での全形式比較によって調べ，
Poissonモジュライの既知の相対幾何を非可換モジュライへ移す。
全ての縮小は，滑らかな可換模型の全体を保つ。

\subsection{中心方程式とJacobian括弧}

$W$ 上のambientと中心商を $A$，$B=A/(\Omega)$ とし，
$F=F^{\mathrm{sh}}=\operatorname{Sym}^{-1}(\Omega)$ と書く。
中心関係の次数は $e_1=6,e_2=4$，$F$ の $z^2$ 係数は $1/2$ である。
次数1の生成元を $x_0,\ldots,x_{n-1},y,z$，次数2では
$x_0,\ldots,x_n,z$ とする。
最低次数の交換子によって方向を
\[
 [x_i,y]=D_i z,\quad D=\sum_iD_i\partial_{x_i}\quad(d=1),
 \qquad
 [x_i,x_j]=p_{ij}z,\quad
 P=\sum_{i<j}p_{ij}\partial_{x_i}\wedge\partial_{x_j}\quad(d=2)
\]
と表す。$D_i,p_{ij}$ は基底の係数であり，相対導分では微分しない。
次数2のrank条件は，係数環の中のPl\"ucker方程式
\[
 Q_{ijk\ell}=p_{ij}p_{k\ell}-p_{ik}p_{j\ell}
                         +p_{i\ell}p_{jk}=0
\]
を意味する。幾何学的ファイバーだけのrank条件には置き換えない。

\begin{lemma}[任意基底上のJacobian Poisson構造]
\label{lem:shadow-jacobian-poisson}
任意の可換 $\mathbb C$-代数 $R$ 上の重み $e_d$ の多項式 $F$ に対し，
\begin{equation}\label{eq:poisson-shadow-bivectors}
 \Pi_{F,D}=\iota_{dF}(D\wedge\partial_y\wedge\partial_z)
 \quad(d=1),\qquad
 \Pi_{F,P}=\iota_{dF}(P\wedge\partial_z)\quad(d=2)
\end{equation}
は重み0の相対Poisson二重導分であり，$F$ はCasimirである。
従って可換超曲面環とその相対ProjにPoisson構造を誘導する。
構成は任意の基底変更に可換し，$R$ の被約性を要求しない。
\end{lemma}
\begin{proof}
縮約は
$\iota_\alpha(D_1\wedge D_2\wedge D_3)
=\alpha(D_1)D_2\wedge D_3-\alpha(D_2)D_1\wedge D_3
+\alpha(D_3)D_1\wedge D_2$ と定める。
次数1の生成元の括弧は
\begin{equation}\label{eq:degree-one-poisson-shadow-brackets}
 \{x_i,x_j\}=0,\quad \{x_i,y\}=D_iF_z,\quad
 \{z,x_i\}=D_iF_y,\quad \{y,z\}=D(F).
\end{equation}
$D,\partial_y,\partial_z$ は互いに可換する。
$D_0$ が可逆なら $u=x_0/D_0$，$s_i=x_i-(D_i/D_0)x_0$ により
$D=\partial_u$ とでき，中央係数 $R[s_i]$ 上の三変数Jacobian括弧となる。
Jacobi恒等式は二階偏微分の対ごとの相殺で従う。
$D_i$ と $F$ の全係数を独立にした普遍多項式環は $D_0$ の局所化に単射なので，
Jacobiatorの各係数は普遍的に零である。
任意の $R$ への基底変更により，零・冪零方向も含めて結論を得る。

次数2の括弧は
\begin{equation}\label{eq:degree-two-poisson-shadow-brackets}
 \{x_i,x_j\}=p_{ij}F_z,\qquad
 \{x_i,z\}=-\sum_jp_{ij}F_{x_j}.
\end{equation}
$F_i=F_{x_i}$，$a=F_z$ と置き，
$J(f,g,h)=\{f,\{g,h\}\}+\{g,\{h,f\}\}+\{h,\{f,g\}\}$ とする。
Leibniz則と二階偏微分の対称性から
\begin{align}
 J(x_i,x_j,x_k)&=\sum_\ell Q_{ijk\ell}
       (aF_{zx_\ell}-F_\ell F_{zz}),
       \label{eq:poisson-shadow-jacobi-xxx}\\
 J(x_i,x_j,z)&=-\sum_{\ell,m}Q_{ij\ell m}F_\ell F_{zx_m}
       \label{eq:poisson-shadow-jacobi-xxz}
\end{align}
を得る。第一式では各 $\ell$ の係数が
$p_{jk}p_{i\ell}+p_{ki}p_{j\ell}+p_{ij}p_{k\ell}=Q_{ijk\ell}$ となり，
第二式では $aF_{x_\ell x_m}$ の項が相殺する。
従ってPl\"ucker方程式から全生成元のJacobiatorは零であり，
各変数に関する導分性によって全多項式について成立する。

両次数で $\{F,-\}=0$ は三重交代積に $dF$ を二度代入することから従う。
重みは次数1で $6-1-2-3=0$，次数2で $4-1-1-2=0$ である。
従って $(F)$ はPoissonイデアルで，斉次局所化の次数零部分も括弧で閉じる。
その標準開上の括弧を貼り合わせればProjに降下する。
全ては多項式恒等式と局所化によるので，任意の基底変更と両立する。
\end{proof}

\begin{lemma}[対称化した中心方程式の偶性]
\label{lem:shadow-even-z}
現在の $W$ 上では
\[
 F^{\mathrm{sh}}=z^2/2+f(x,y)\quad(d=1),\qquad
 F^{\mathrm{sh}}=z^2/2+f(x)\quad(d=2)
\]
と書ける。従って最低次数のPoisson括弧は，
それぞれ $\{x_i,y\}=D_i z$，$\{x_i,x_j\}=p_{ij}z$ である。
\end{lemma}
\begin{proof}
両ambientの表示関係式
\eqref{eq:simultaneous-degree-one-relations}，
\eqref{eq:simultaneous-ambient-relations}は，
$x_i$ と $y$ を固定し $z\mapsto-z$ とする基底線型の反自己同型 $\rho$ を持つ。
最低交換子の関係式はその負へ移り，$z$ を含む残りの交換子の関係式は自身へ移る。
完全対称化された語の順序反転は置換の集合を保つ。
非零方向上では補題\ref{lem:normalized-ambient-center-near-zero}から
\[
 Z(A)_{e_d}=\operatorname{Sym}^{e_d}K\oplus\mathcal O_W\Omega
\]
である。$\rho\Omega$ の $z^2$ 係数も $1/2$ なので，
$\rho\Omega-\Omega=p\in\operatorname{Sym}^{e_d}K$ となる。
$\rho$ は $p$ を固定するから，$\rho^2=1$ より $2p=0$，従って $p=0$ である。
$W$ は被約，非零方向は稠密，$A_{e_d}$ は局所自由なので，
$\rho\Omega=\Omega$ は $W$ 全体で成立し，任意基底変更後にも保たれる。
重みと $z^2$ の係数から結論を得る。
\end{proof}

\subsection{全同値への降下と零切断の検出}

\begin{proposition}[Poisson模型への表現可能な射]
\label{prop:poisson-shadow-morphism}
$\mathcal P_{n,d}$ を滑らかな次数 $d$ のdel Pezzo族と相対Poisson構造の
モジュライとする。式\eqref{eq:poisson-shadow-bivectors}は
\[
 \Theta:\mathfrak U_{n,d}\longrightarrow\mathcal P_{n,d},
 \qquad [B]\longmapsto(X^{\mathrm{sh}},\Pi^{\mathrm{sh}})
\]
を定める。Poisson構造を忘れた合成は既存の $\operatorname{sh}$ であり，
$\Theta$ は代数空間で表現可能である。
Poisson零切断の逆像は，可換零切断 $\mathcal D_{n,d}$ とschemeとして一致し，
その制限は恒等である。
\end{proposition}
\begin{proof}
まず全同値に対する自然性を確認する。
fppf局所的に同値を通常の次数付き同型 $\phi:B_a\to B_b$ として表す。
補題\ref{lem:all-graded-isom-weighted-linear}により，次数1では
\[
 \phi(x)=Mx',\quad\phi(y)=\beta y',\quad\phi(z)=\gamma z',
 \qquad \beta M D_b=\gamma D_a,
\]
次数2では
\[
 \phi(x)=Mx',\quad\phi(z)=cz',\qquad MP_bM^{\mathsf t}=cP_a
\]
となる。方向の等式は最低交換子と $z'$ の自由直和因子から得られ，
方向座標を消去する操作は使わない。
同型はambientへ持ち上がり，中心関係を単元倍へ送る。
完全対称化と $z^2$ 係数の比較から，正確に
\[
 \phi(F_a)=\gamma^2F_b\quad(d=1),\qquad
 \phi(F_a)=c^2F_b\quad(d=2)
\]
である。

導分の移送を $\phi_*\delta=\phi\delta\phi^{-1}$ とすると
\[
 \phi_*(D_a\wedge\partial_y\wedge\partial_z)
   =\gamma^{-2}D_b\wedge\partial_{y'}\wedge\partial_{z'},
 \qquad
 \phi_*(P_a\wedge\partial_z)=c^{-2}P_b\wedge\partial_{z'}.
\]
これらの係数と方程式の係数が相殺し，
$\phi_*\Pi_a^{\mathrm{sh}}=\Pi_b^{\mathrm{sh}}$ となる。
相対微分では基底の単元を微分しない。
従ってPoisson構造は単元倍まででなく厳密に保たれる。
次数スカラーはProj上恒等であり，全ての式は合成・逆射・基底変更と両立する。
全同値群oidに沿う忠実平坦降下から $\Theta$ を得る。
その相対自己同型の核は $\operatorname{sh}$ の相対自己同型の核に含まれるため自明であり，
命題\ref{prop:representable-commutative-shadow}と同じ判定から表現可能性を得る。

零切断のscheme構造を確認するため，任意の係数環 $R$ 上で考える。
次数1の $D_+(x_0)$ 上で $u_i=x_i/x_0$，$v=y/x_0^2$，$w=z/x_0^3$ と置くと
\[
 \{u_i,v\}=(D_i-D_0u_i)w.
\]
このchartの環はmonicな $w^2$ 関係により $R[u_i,v]$ 上で基底 $1,w$ を持つ。
従って括弧が零なら $D_i-D_0u_i=0$ であり，$n\ge2$ から全ての $D_i=0$ を得る。
次数2で $u_i=x_i/x_0$，$v=z/x_0^2$ と置くと
\[
 \{u_i,u_j\}=v(p_{ij}-p_{i0}u_j-p_{0j}u_i).
\]
同じmonic分裂と $n\ge2$ より，全ての $p_{ij}=0$ を得る。
逆は構成式から明らかである。
この同値は任意の非被約 $R$ でも成立するため，零部分関手が一致し，
$W$ 上の方向イデアルそのものが検出される。
零切断では $\Omega=F$ が元の可換方程式なので，その上の射は恒等である。
\end{proof}

\begin{remark}[規格化と中心選択]
同じ中心イデアルの生成元を単元倍しても，その対称化 $F$ を
$F/(2[z^2]F)$ と規格化すれば上のPoisson構造は変わらない。
非零方向で異なる規格化中心元の差が
$p\in\operatorname{Sym}^{e_d}K$ なら，$p$ は三重導分の各因子で消えるので
$\Pi_{F+p}=\Pi_F$ はambient上で成立する。
ただし $F=0$ と $F+p=0$ は一般に異なる中心商である。
ここで構成した射は実際の中心商の関係式に付随する。
次数2では第\ref{sec:degree-two-shadow-degree}節で，規格化中心のtorsorを用い，
異なる全方向中心選択が共通の一般開上で同じ比較射を与えることを示す。
選んだ有限型開領域全体の一致や既存のformality量子化との全次数での一致は主張しない。
\end{remark}

\begin{proposition}[可換点での全形式比較]
\label{prop:poisson-shadow-formal-comparison}
任意の滑らかな可換点 $X$ で，$\Theta$ はframed形式変形群oidの同値を誘導する。
接写像は，通常変形と一次交換子によるPoisson方向の同定の下で恒等である。
\end{proposition}
\begin{proof}
方向イデアルを $J$ とする。
定理\ref{thm:simultaneous-degree-one-central-family}，
\ref{thm:simultaneous-degree-two-central-family}の規格化から，
両次数で $F^{\mathrm{sh}}\equiv F\pmod{J^2}$ である。
式\eqref{eq:poisson-shadow-bivectors}は方向に線型なので，
$\Pi^{\mathrm{sh}}-\Pi_F\in J^3$ となる。
従って可換係数の方向では元の可換変形を返し，
Poisson方向では有限中心商の一次交換子と同じ括弧を返す。
商と局所化のLeibniz則も一致するから，これはtailsの初交換子にも一致する
（定理\ref{thm:simultaneous-central-all-artin}の証明参照）。
標準方程式を $2F$ に直す規約ではJacobianパラメータを半分にすればよく，
この一定の規格化を双方に施して比較する。

可換係数の横断sliceを選ぶと，sourceとtargetのframed形式関手は
定理\ref{thm:abel-formal-ring}，\ref{thm:simultaneous-central-all-artin}および
Poisson構造定理によって同じ完備Noether局所環 $R_X$ で表される。
したがって $\Theta$ は余接空間上可逆な $R_X$ の自己準同型に対応し，
補題\ref{lem:formal-endomorphism-inverse}から同型である。
次数2のPl\"ucker特異性も含めて全形式環を比較している。
framed自己同型は両側で自明なので，同型類の自然な全単射は群oid同値を与える。
\end{proof}

\subsection{Poisson shadowの\'{e}tale性と平坦な可換模型}

以下では，命題\ref{prop:poisson-shadow-morphism}で構成した
全同値と任意基底変更に両立する射
\[
 \Theta:\mathfrak U_{n,d}\longrightarrow\mathcal P_{n,d},
 \qquad p\Theta=\operatorname{sh}
\]
を用いる。可換零切断を $j:\mathcal D_{n,d}\hookrightarrow\mathfrak U_{n,d}$，
Poisson零切断を $s:\mathcal D_{n,d}\hookrightarrow\mathcal P_{n,d}$ と書く。
入力は $\Theta j=s$ と，命題\ref{prop:poisson-shadow-formal-comparison}の
可換点でのframed完成同型，および
命題\ref{prop:poisson-shadow-morphism}のスキーム論的な零切断の一致である。

\begin{theorem}[零切断を保つPoissonモジュライへの\'{e}tale射]
\label{thm:poisson-shadow-etale}
全可換零切断を保って $\mathfrak U_{n,d}$ をさらに開部分へ制限すると，
$\Theta$ はrepresentable，separated，有限表示かつ\'{e}taleとなり，
全点でstabilizer preservingである。さらに
\[
 \mathfrak U_{n,d}\times_{\mathcal P_{n,d}}\mathcal D_{n,d}
 \simeq\mathcal D_{n,d}
\]
であり，この同型は零切断を与える。
$\Theta$ の像はPoisson零切断全体を含む開部分スタックである。
\end{theorem}
\begin{proof}
\emph{表現可能性。}
$\Theta$ のrepresentabilityは命題\ref{prop:poisson-shadow-morphism}による。
始域・終域は有限型の複素代数スタックで，終域の対角は有限であるため，
$\Theta$ は有限表示である。
また始域は分離的で，終域の対角は分離的なので $\Theta$ はseparatedである。

\emph{可換点での形式同型。}
命題\ref{prop:poisson-shadow-formal-comparison}により，
可換模型 $X$ における $\Theta$ のframed完成は同型である。
可換点の自己同型群への写像も，元の可換模型を保つ恒等写像である。

終域の\'{e}tale表示で基底変更し，始域も\'{e}tale局所的にスキームで表す。
選んだ複素点での完備局所環写像は上のframed形式同型となるから，
\cite[Lemma 41.11.3, Tag 039M]{StacksProject}によりそこで\'{e}taleである。
よって $\Theta$ の\'{e}tale locusは全可換零切断を含む開部分スタックである。
有限型複素表示上では，零切断の全閉点を含むことから零切断全体を含むと分かる。
この開集合へ制限する。

\emph{全点での安定化群。}
$\mathcal P_{n,d}$ の対角は
\cite[定理9.2]{PoissonDP}により有限なので，その慣性はproperである。
\cite[Lemma 101.45.4, Tag 0DUA]{StacksProject}によって
$\Theta$ のstabilizer preserving locusは開である。
零切断では既に全自己同型群に同型を誘導するから，
この開部分への制限も全零切断を保つ。

前項の零イデアルの計算から，零切断の逆像は既に所望のものに等しい。
最後に\'{e}tale射は開射であり，像は $\Theta j=s$ により全Poisson零切断を含む。
\end{proof}

\begin{corollary}[可換模型への平坦射と相対特異性]
\label{cor:shadow-flat-cone-fibers}
定理\ref{thm:poisson-shadow-etale}の近傍では
\[
 \operatorname{sh}:\mathfrak U_{n,d}\longrightarrow\mathcal D_{n,d}
\]
はquasi-affine，有限表示，平坦かつ全射である。
純相対次元は $c_{n,1}=n$，$c_{n,2}=2n-1$ である。
任意の幾何学的点 $[X]:\operatorname{Spec}K\to\mathcal D_{n,d}$ に沿うファイバーは，
相対的な自己同型も取り除いたスキームであり，Jacobi錐
\[
 C_{n,1,K}=\mathbb A_K^n,\qquad
 C_{n,2,K}=\operatorname{Cone}_K\operatorname{Gr}(2,n+1)
\]
へのseparatedな\'{e}tale射を持つ。その像は頂点を含む開部分であり，
頂点の逆像は可換模型を表す一点である。
各幾何学的ファイバーは正規，Cohen--MacaulayかつGorensteinである。

次数1と次数2の $n=2$ では $\operatorname{sh}$ は滑らかである。
次数2の $n\ge3$ では，その相対非smooth locusの台はちょうど可換零切断である。
零点での相対特異性は $n=3$ なら五次元の非退化二次錐，
$n\ge4$ なら非局所完全交叉であり，相対接空間の次元超過は
\[
 \binom{n+1}{2}-(2n-1)=\frac{(n-1)(n-2)}2
\]
である。
\end{corollary}
\begin{proof}
入力文献\cite[定理8.8]{PoissonDP}によれば，
$p:\mathcal P_{n,d}\to\mathcal D_{n,d}$ はaffine，有限表示かつ平坦であり，
幾何学的ファイバーは表示した錐である。
その前に\'{e}tale射 $\Theta$ を合成すれば，平坦性・有限表示性を得る。
合成は零切断 $j$ を持つので全射である。

$\Theta$ はseparatedで有限表示の\'{e}tale射なのでquasi-finiteである。
代数空間版Zariski主定理\cite[Tag 082J]{StacksProject}によりquasi-affineであり，
affine射 $p$ との合成もquasi-affineである。
したがってスキームで表される各古典族に沿う基底変更もスキームとなる。
上記の幾何学的ファイバーについては
\[
 \Theta_X:(\mathfrak U_{n,d})_X\longrightarrow C_{n,d,K}
\]
が\'{e}taleであり，零切断の逆像に関する前定理を基底変更できる。
次元，正規性，Cohen--Macaulay性，Gorenstein性は
この\'{e}tale写像と錐の対応する性質から従う。
Pl\"ucker錐のGorenstein性は
\cite[命題10.4，\texttt{prop:plucker-terminal}]{PoissonDP}で
標準加群 $\omega_{\mathbb C[C]}\simeq\mathbb C[C](-(n+1))$ として計算されている。
$n=2$ では錐は単に三次元affine空間である。
各非空成分は純次元 $c_{n,d}$ の錐の開部分へ\'{e}taleに写るので同じ純次元を持つ。

滑らかさと非smooth locusの台は $p$ についての対応する性質を
\'{e}taleに引き戻せばよい。
零切断の逆像の一致により，零点の余分な分枝は生じない。
局所完全交叉性，二次錐の形，接空間次元は
\cite[補題10.1，\texttt{lemma:cone-lci}]{PoissonDP}と
\'{e}tale局所環の性質から従う。
\end{proof}

\begin{remark}[ファイバーの範囲]
上のファイバーは，$X$ との同型を記憶する
$\operatorname{Spec}K\to\mathcal D_{n,d}$ に沿うファイバーである。
それを $C_{n,d,K}$ 全体やその一つの開部分と同一視したのではなく，
その開部分への\'{e}tale射を得ている。
ファイバー全体の既約性・連結性や，$\Theta_X$ の次数1は結論しない。
また，自己同型も割る粗shadowのファイバーとは異なる。
\end{remark}

\subsection{粗Poisson空間との比較と粗shadowの平坦領域}

古典・Poisson・非可換モジュライの粗空間をそれぞれ
$M^{\mathrm{cl}}_{n,d}$，$N_{n,d}$，
$M^{\mathrm{nc}}_{n,d}$ と書く。
前二者は入力文献\cite[定理9.2]{PoissonDP}の $M_{n,d},N_{n,d}$ である。

\begin{theorem}[粗shadowの\'{e}tale表示と平坦領域]
\label{thm:coarse-shadow-etale-flat-locus}
定理\ref{thm:poisson-shadow-etale}の近傍では，Cartesian図式
\[
\begin{array}{ccc}
\mathfrak U_{n,d}&\xrightarrow{\ \Theta\ }&\mathcal P_{n,d}\\
\downarrow&&\downarrow\\
M^{\mathrm{nc}}_{n,d}&\xrightarrow{\ \overline\Theta\ }&N_{n,d}
\end{array}
\]
があり，$\overline\Theta$ はseparated，有限表示かつ\'{e}taleである。
$M^{\mathrm{nc}}_{n,d}$ は正規Cohen--Macaulayな整quasi-affineスキームである。
粗shadow
\[
 \overline{\operatorname{sh}}:
 M^{\mathrm{nc}}_{n,d}\longrightarrow M^{\mathrm{cl}}_{n,d}
\]
は全射であり，その平坦領域は厳密に
\[
 \operatorname{Flat}(\overline{\operatorname{sh}})
 =
 \overline{\operatorname{sh}}^{-1}
 \bigl((M^{\mathrm{cl}}_{n,d})^{\mathrm{sm}}\bigr)
\]
となる。
したがって古典粗空間の特異点上では，
構成近傍に属するファイバーのどの点でも粗shadowは平坦でない。
\end{theorem}
\begin{proof}
$\Theta$ はseparated，\'{e}taleかつstabilizer preservingである。
系\ref{cor:algebraic-coarse-local-model}の証明と同じく，
\cite[Tags 0DUS, 0DUE]{StacksProject}を
well-nigh affineな\'{e}tale被覆上で適用して降下すれば，
粗空間上のCartesian図式と $\overline\Theta$ の\'{e}tale性を得る。
分離性と有限表示性も成立する。
$N_{n,d}$ は
\cite[定理9.2]{PoissonDP}によってaffineスキームなので，
quasi-finiteかつseparatedな $\overline\Theta$ に
\cite[Tag 082J]{StacksProject}を適用すると，始域はquasi-affineスキームである。
Poisson粗空間の正規性は正規な座標環の不変部分環から，
Cohen--Macaulay性は\cite[定理9.6]{PoissonDP}から従う。
これらは\'{e}taleに引き戻せる。
整性は既存の $M^{\mathrm{nc}}_{n,d}$ の整性である。
粗shadowの全射性は全可換模型の包含から従う。

Poisson粗忘却射を $\overline p:N_{n,d}\to M^{\mathrm{cl}}_{n,d}$ とする。
固定commitの\cite[定理9.6，\texttt{thm:coarse-flat-locus}]{PoissonDP}は
\[
 \operatorname{Flat}(\overline p)
 =\overline p^{-1}\bigl((M^{\mathrm{cl}}_{n,d})^{\mathrm{sm}}\bigr)
\]
を与える。点 $u\mapsto v\mapsto x$ における局所環で，
$\mathcal O_{N,v}\to\mathcal O_{M^{\mathrm{nc}},u}$ は
\'{e}taleな局所準同型，従って忠実平坦である。
平坦性の合成と忠実平坦降下により，
$\overline p$ が $v$ で平坦であることと，
$\overline{\operatorname{sh}}=\overline p\,\overline\Theta$ が $u$ で平坦であることは同値である。
これで記載した平坦領域の等式を得る。
\end{proof}

\begin{remark}[粗ファイバーに残る冪零構造]
スタックのshadowの幾何学的ファイバーは正規Gorensteinであるが，
粗shadowの幾何学的ファイバーには冪零構造が生じ得る。
入力文献\cite[定理9.6]{PoissonDP}が与える
$(\overline p^{-1}([X]))_{\mathrm{red}}\simeq
C_{n,d}/\operatorname{Aut}(X)$ は被約化についての等式である。
本稿の粗ファイバーはこのPoisson粗ファイバーへの\'{e}tale射を持つ。
したがってその被約化は上記の有限群商への\'{e}tale射を持つが，
ファイバー自体の被約性や全体の既約性は主張しない。
\end{remark}

\subsection{零切断のblow-upと相対Pl\"ucker特異性}

\begin{corollary}[零切断を中心とする滑らかな解消]
\label{cor:nc-shadow-zero-blowup-resolution}
次数2において
\[
 \rho:\widetilde{\mathfrak U}_{n,2}
 =\operatorname{Bl}_{\mathcal D_{n,2}}\mathfrak U_{n,2}
 \longrightarrow\mathfrak U_{n,2}
\]
と置く。この射はrepresentable，射影的かつ固有で，零切断の外では同型である。
$\widetilde{\mathfrak U}_{n,2}$ は滑らかなDMスタックであり，
$\widetilde{\mathfrak U}_{n,2}\to\mathcal D_{n,2}$ も滑らかである。
従って $n\ge3$ では $\rho$ は特異性の解消である。

例外因子 $E$ はPoisson側の相対Grassmann族と標準的に同型であり，
各古典模型上のファイバーは $\operatorname{Gr}(2,n+1)$ である。
その法線束は，相対Poissonベクトル束の射影化の中の
tautological line bundle $\mathcal O(-1)$ である。
さらに $\mathfrak U_{n,2}$ はGorensteinで，滑らかなatlasの意味でterminalであり，
\[
 K_{\widetilde{\mathfrak U}_{n,2}}
 =\rho^*K_{\mathfrak U_{n,2}}+nE
\]
が成り立つ。これは粗空間のGorenstein性・terminal性の主張ではない。
次数1の $\mathfrak U_{n,1}$ は滑らかである。
\end{corollary}
\begin{proof}
$\Theta$ は平坦で，零切断のスキーム論的逆像は可換零切断に等しい。
blow-upと平坦基底変更の可換性\cite[Tag 0805]{StacksProject}により
\[
 \widetilde{\mathfrak U}_{n,2}
 \simeq
 \mathfrak U_{n,2}\times_{\mathcal P_{n,2}}
 \operatorname{Bl}_{\mathcal D_{n,2}}\mathcal P_{n,2}.
\]
右辺からPoisson側のblow-upへの射は\'{e}taleである。
\cite[定理10.3，\texttt{thm:poisson-resolution}]{PoissonDP}によって
後者は相対Grassmann族上の $\mathcal O(-1)$ の全空間であり，
古典スタック上滑らかである。
これで絶対・相対の滑らかさ，射影性と零切断外の同型性を得る。

Poisson側の例外因子は零切断を経由して写るので，
その基底変更には
$\mathfrak U_{n,2}\times_{\mathcal P_{n,2}}\mathcal D_{n,2}
=\mathcal D_{n,2}$ を使える。
したがって例外因子はPoisson側のもの自体に同型であり，法線束も引き戻される。
錐のGorenstein terminal性とdiscrepancy $n$ は
\cite[命題10.4，\texttt{prop:plucker-terminal}]{PoissonDP}による。
滑らかな古典係数方向との積，および\'{e}tale写像でこれらの性質を移せる。
$n=2$ の場合は三次元affine空間の原点のblow-upであり，
同じdiscrepancy $2=n$ を持つ。
次数1についてはPoisson側が滑らかであることを用いる。
\end{proof}

\subsection{開埋込みへの残る障害}

\begin{example}[形式恒等な\'{e}tale射が持つ別の分枝]
\label{ex:etale-shadow-not-open-immersion}
$k=\mathbb C$，
\[
 V=\operatorname{Spec}k[t,(1+t^2)^{-1},(1+3t^2)^{-1}]
\]
とし，$\mu_2$ は $t$ に符号で作用する。
$h(t)=t+t^3$ は符号作用と両立し，
\[
 H:[V/\mu_2]\longrightarrow[\mathbb A^1_s/\mu_2]
\]
を与える。両辺は有限対角を持つ整な滑らかな有限型DMスタックであり，
$H$ はrepresentable，separatedかつ\'{e}taleで，全点でstabilizer preservingである。
零gerbe $B\mu_2$ 上では恒等であり，そのスキーム論的逆像も同じ零gerbeである。
零点の形式環の射 $s\mapsto t+t^3$ は同型で，接写像は恒等である。
それでも，零gerbeを含む始域のどのZariski開部分へ制限しても，
$H$ は開埋込みにならない。
\end{example}
\begin{proof}
$h'(t)=1+3t^2$ は $V$ 上可逆なので $h$ は\'{e}taleである。
同じ群による同変射であり，標的のatlas $\mathbb A^1_s$ による基底変更は
$h:V\to\mathbb A^1_s$ だから，商の射もrepresentableかつ\'{e}taleである。
$1+t^2$ を可逆化したので $h^{-1}(0)$ はスキーム論的にも $t=0$ だけである。
非零点では両側の安定化群は自明，零点ではどちらも $\mu_2$ である。
形式同型は線型係数が $1$ であることから従う。

粗空間の座標を $a=t^2$，$b=s^2$ とすると，
粗射は
\[
 b=a(1+a)^2
\]
であり，関数体拡大の次数は $3$ である。
非空開部分への制限では関数体は変わらないので，
どの零保持開制限も開埋込みとはならない。
実際，schemeの関係 $V\times_{\mathbb A^1}V$ の対角以外の部分は
$1+t^2+tu+u^2=0$ で与えられる。
そのsource像は零点を含まないが，稠密であり閉包は零点を含む。
対角が開閉であることや形式近傍が一つであることだけでは，
この像を零保持開集合から除去できない。
\end{proof}

\begin{proposition}[大域比較を決める一般次数]
\label{prop:poisson-shadow-degree-one-criterion}
定理\ref{thm:coarse-shadow-etale-flat-locus}の射に対して，
\[
 \delta_{n,d}
 =
 [\,K(M^{\mathrm{nc}}_{n,d}):K(N_{n,d})\,]
\]
は有限の正整数である。次の条件は同値である。
\begin{enumerate}
\item $\delta_{n,d}=1$。
\item $\overline\Theta$ は開埋込みである。
\item $\Theta$ は開埋込みである。
\end{enumerate}
この一般次数は非空開部分への縮小で変わらない。
本節の形式比較と\'{e}tale性だけでは $\delta_{n,d}=1$ は結論できない。
\end{proposition}
\begin{proof}
両粗空間は整であり，$\overline\Theta$ は非空な開像を持つのでdominantである。
有限型かつquasi-finiteであることから，関数体拡大は有限である。
次数が1なら，正規な $N_{n,d}$ への
birational，quasi-finiteかつseparatedな射となる。
\cite[Tag 082J]{StacksProject}のZariski主定理で，
終域のその関数体内での正規化は終域自身だから，$\overline\Theta$ は開埋込みである。
Cartesian図式によりこれは $\Theta$ の開埋込みを与える。
逆に $\Theta$ が開埋込みなら，粗空間の開部分への制限の性質によって
$\overline\Theta$ も開埋込みであり，関数体は等しい。
非空開制限で関数体は不変である。
前例では全ての局所条件が成立する一方で $\delta=3$ である。
\end{proof}

\begin{remark}
全零切断での形式同型は，対角の近傍にある一意な枝を同定する。
大域の開埋込みには，それに加えて離れた同一Poisson模型の枝を排除する必要がある。
本稿の具体的な $\Theta$ について，
定理\ref{thm:degree-two-shadow-degree-four}と\ref{thm:degree-one-shadow-degree-five}により
$\delta_{n,2}=4$，$\delta_{n,1}=5$ となり，両次数で開埋込みでない。
両次数の開像と補集合，および逆量子化のさらなる幾何的記述は大域比較の課題として残る。
\end{remark}
